% ===== Document Preamble =====
\documentclass[11pt]{article}
\usepackage[margin=0.5in]{geometry}
\usepackage{amsmath,amssymb}
\usepackage{xcolor}
\usepackage{enumitem}
\usepackage{needspace}
\usepackage{comment}

% ===== Colors =====
\definecolor{AnsRed}{RGB}{170,0,0}
\definecolor{WrongBlue}{RGB}{0,0,160}

% ===== Choice list styling (tight, consistent) =====
\setlist[enumerate]{%
  leftmargin=1.5em,
  itemsep=0.25em,
  topsep=0.25em,
  label=\textbf{\Alph*.}
}

% ===== Style Wrappers =====
\newcommand{\QuestionStyle}[1]{{\large\bfseries\boldmath #1}}
\newcommand{\AnswerStyle}[1]{{\large\bfseries\boldmath #1}}
\newcommand{\ReasonStyle}[1]{{\normalsize #1}}

% ===== Show / Hide Reasoning Toggle =====
\newif\ifshowreason
\showreasontrue      % Instructor version
% \showreasonfalse   % Student version

% ===== Choice Item Macros (Answer + Reasoning) =====
% Usage:
%   \Correct{Answer}{Reasoning}
%   \Wrong{Answer}{Reasoning}

\newcommand{\Correct}[2]{%
  \item \AnswerStyle{\textcolor{AnsRed}{#1}}%
  \ifshowreason\\[-0.2em]
    \ReasonStyle{\textcolor{AnsRed}{#2}}%
  \fi
}

\newcommand{\Wrong}[2]{%
  \item \AnswerStyle{\textcolor{WrongBlue}{#1}}%
  \ifshowreason\\[-0.2em]
    \ReasonStyle{\textcolor{WrongBlue}{#2}}%
  \fi
}

% ===== Question Block Macro =====
% Usage: \QBlock{Q1}{stem...}{<items>}

\newcommand{\QBlock}[3]{%
  \Needspace{8\baselineskip}%
  \noindent
  \begin{minipage}{\linewidth}
    \QuestionStyle{#1.~#2}

    \vspace{0.25em}
    \begin{enumerate}
      #3
    \end{enumerate}

    \vspace{0.25em}
    \hrule height 0.4pt
  \end{minipage}
  \vspace{0.75em}
}

% ===== Document Start =====

\begin{document}

\begin{center}
{\huge \textbf{Physics Isomorphic Questions}}\\[0.25em]
{\Large Modeling, Units, Vectors, and Mechanics}\\[0.25em]
{\large Multiple choice questions generated below. Answers in RED and incorrect answers in BLUE.}
\end{center}

\vspace{0.75em}

\QBlock{Q1}{Which of the following modeling statements is most appropriate?}{
  \Wrong{To compute the torque on a door about its hinges when you push near the handle, it is safe to model the door as a point mass located at its center.}{This treats a rotational/rigid-body problem as pure translation. Torque depends on the lever arm and force application point, so a point-mass model throws away the geometry needed for rotation.}
  \Correct{To estimate how long a city bus takes to travel between two stops that are $8\ \mathrm{km}$ apart, it is usually reasonable to model the bus as a point mass when tracking its position along the route.}{When the distance scale of interest ($\sim 8\ \mathrm{km}$) is much larger than the bus's size ($\sim 10\ \mathrm{m}$), its internal dimensions are negligible for center-of-mass motion. In that regime, treating it as a point mass is a standard approximation for modeling translation (position vs.\ time) without resolving rotations or shape.}
  \Wrong{Even if you decide to ignore air resistance on a thrown ball, you must still include detailed airflow effects because the ball is not literally in a vacuum.}{This confuses a physical idealization with a literal claim about the environment. ``Ignore drag'' is equivalent to using a vacuum-force model for the motion equation.}
  \Wrong{None of the above statements is appropriate.}{Option B is a valid scale-separation approximation for translational motion over kilometers. Therefore it is not correct that none are appropriate.}
}

\QBlock{Q2}{Which of the following unit-conversion equations is correct?}{
  \Wrong{$1\ \mathrm{mm^2}=10^{-3}\ \mathrm{m^2}$}{This applies the $10^{-3}$ factor for millimeters to area without squaring it. Area scales with the square of the length conversion.}
  \Correct{$1\ \mathrm{mm^2}=10^{-6}\ \mathrm{m^2}$}{Since $1\ \mathrm{mm}=10^{-3}\ \mathrm{m}$, squaring both sides gives $1\ \mathrm{mm^2}=(10^{-3}\ \mathrm{m})^2=10^{-6}\ \mathrm{m^2}$. The key idea is that powers apply to the conversion factor as well.}
  \Wrong{$1\ \mathrm{mm^2}=10^{6}\ \mathrm{m^2}$}{This inverts the metric prefix relationship. Millimeters are smaller than meters, so $\mathrm{mm^2}$ must also be smaller than $\mathrm{m^2}$.}
  \Wrong{$1\ \mathrm{mm^2}=10^{-2}\ \mathrm{m^2}$}{This treats ``milli-'' like ``centi-'' (a factor of $10^{-2}$) and also ignores squaring. The correct factor is $10^{-6}$.}
}

\QBlock{Q3}{Which option has the incorrect SI unit?}{
  \Wrong{Pressure: $\mathrm{N/m^2}$}{Pressure is force per area, so $\mathrm{Pa}=\mathrm{N/m^2}$ is correct.}
  \Wrong{Angular momentum: $\mathrm{kg\,m^2/s}$}{Angular momentum has units $\mathrm{(mass)(length^2)/(time)}$, so $\mathrm{kg\,m^2/s}$ is correct.}
  \Wrong{Electric field: $\mathrm{N/C}$}{Electric field is force per charge, so $\mathrm{N/C}$ is correct.}
  \Correct{Power: $\mathrm{J\cdot s}$}{Power is energy per unit time, so its SI unit is $\mathrm{W}=\mathrm{J/s}$. Writing $\mathrm{J\cdot s}$ multiplies by time instead of dividing, which corresponds to a different physical dimension.}
}

\QBlock{Q4}{A search-and-rescue drone's position relative to the command tent (at the origin) is $\vec r = 4\,\hat\imath + 11\,\hat\jmath$ in kilometers. Which statement is correct?}{
  \Wrong{The drone is $15\ \mathrm{km}$ from the command tent.}{This incorrectly adds components as $4+11$. The distance is the magnitude $|\vec r|=\sqrt{x^2+y^2}$, not the Manhattan sum.}
  \Wrong{The drone is $11\ \mathrm{km}$ from the command tent.}{This confuses the $y$-component with the total distance. The drone also has an $x$-offset.}
  \Wrong{The drone is $4\ \mathrm{km}$ from the command tent.}{This confuses the $x$-component with the total distance. Both components contribute to the magnitude.}
  \Correct{None of the above is correct.}{The distance from the origin is $|\vec r|=\sqrt{4^2+11^2}=\sqrt{137}\approx 11.70\ \mathrm{km}$. Since $11.70$ km is not equal to $15$, $11$, or $4$, none of A--C is correct.}
}

\QBlock{Q5}{A cart's 1D displacement is given by $\vec s(t)=\big((1.10\,\mathrm{m/s^3})t^3-(2.40\,\mathrm{m/s^2})t^2\big)\,\hat\imath$. Which statement is correct?}{
  \Wrong{This motion has constant acceleration.}{A $t^3$ term in $s(t)$ produces a time-dependent acceleration after two derivatives. Constant acceleration would require $s(t)$ to be at most quadratic in $t$.}
  \Correct{The magnitude of the acceleration at $t=0$ is $4.8\ \mathrm{m/s^2}$.}{Compute velocity and acceleration: $v(t)=\frac{ds}{dt}=3(1.10)t^2-2(2.40)t=3.30t^2-4.80t$ and $a(t)=\frac{dv}{dt}=6.60t-4.80$. Thus $a(0)=-4.80\ \mathrm{m/s^2}$, so the magnitude is $|a(0)|=4.8\ \mathrm{m/s^2}$.}
  \Wrong{The magnitude of the acceleration is $2.4\ \mathrm{m/s^2}$ for all $t$.}{This mistakes a coefficient in $s(t)$ for the acceleration. Acceleration comes from $\frac{d^2s}{dt^2}$, which here is $6.60t-4.80$, not a constant $2.4$.}
  \Wrong{The magnitude of the acceleration at $t=0$ is $0\ \mathrm{m/s^2}$.}{This incorrectly assumes plugging in $t=0$ forces acceleration to vanish. Here $a(t)$ has a constant term, so $a(0)\neq 0$.}
}

\QBlock{Q6}{A robot moves along a straight path from position A, $\vec r_A=1.5\,\hat\imath-3.0\,\hat\jmath$ (m), to position B, $\vec r_B=-5.5\,\hat\imath+4.5\,\hat\jmath$ (m), taking $4.0\ \mathrm{s}$. Which statement is correct?}{
  \Wrong{The distance traveled is $14.5\ \mathrm{m}$.}{This adds component changes as $|\Delta x|+|\Delta y|=7.0+7.5=14.5$, which is not the straight-line distance. For a straight path, the distance is the magnitude $\sqrt{(\Delta x)^2+(\Delta y)^2}$.}
  \Wrong{The speed at $t=0$ must be $2.6\ \mathrm{m/s}$.}{This treats the average speed as if it uniquely determines the initial speed. With only endpoints and total time, many velocity-time histories are possible.}
  \Correct{The distance traveled is $10.3\ \mathrm{m}$.}{Because the robot travels straight from A to B, the distance equals the displacement magnitude. $\Delta\vec r=( -5.5-1.5)\hat\imath+(4.5-(-3.0))\hat\jmath=-7.0\hat\imath+7.5\hat\jmath$, so $|\Delta\vec r|=\sqrt{(-7.0)^2+(7.5)^2}=\sqrt{105.25}\approx 10.26\ \mathrm{m}\approx 10.3\ \mathrm{m}$.}
  \Wrong{The acceleration at $t=0$ must be $0.64\ \mathrm{m/s^2}$.}{Acceleration cannot be determined from just start/end position and elapsed time. Infinitely many accelerations could connect the same endpoints in $4.0$ s.}
}

\QBlock{Q7}{A frog leaps from level ground with speed $v_0=7.5\ \mathrm{m/s}$ at an angle $\theta=40^\circ$ above the horizontal (ignore air resistance). What is the frog's speed at maximum height?}{
  \Correct{$5.7\ \mathrm{m/s}$}{In ideal projectile motion, the horizontal component stays constant: $v_x=v_0\cos\theta$. At maximum height the vertical component is zero ($v_y=0$), so the speed is $v=\sqrt{v_x^2+0^2}=v_0\cos\theta=7.5\cos 40^\circ\approx 5.75\ \mathrm{m/s}$.}
  \Wrong{$7.5\ \mathrm{m/s}$}{This ignores that the vertical component decreases to zero at the top. Only the horizontal component remains.}
  \Wrong{$4.8\ \mathrm{m/s}$}{This corresponds to $v_0\sin\theta$, the initial vertical component, not the speed at maximum height where $v_y=0$.}
  \Wrong{$0\ \mathrm{m/s}$}{At the top only the vertical velocity is zero. The horizontal velocity generally remains nonzero.}
}

\QBlock{Q8}{A glider on a frictionless air track has initial velocity $\vec v_0=(12.0\ \mathrm{m/s})\,\hat\imath$ at $t=0$. A fan provides a constant acceleration $\vec a=(-1.5\ \mathrm{m/s^2})\,\hat\imath+(2.0\ \mathrm{m/s^2})\,\hat\jmath$. Which statement is correct?}{
  \Wrong{The $x$-component of the glider's velocity remains constant.}{Because $a_x=-1.5\ \mathrm{m/s^2}\neq 0$, the $x$-component changes as $v_x(t)=v_{0x}+a_xt$.}
  \Wrong{The $y$-component of the glider's velocity remains constant.}{Because $a_y=2.0\ \mathrm{m/s^2}\neq 0$, the $y$-component changes from its initial value ($0$) to nonzero values.}
  \Wrong{The speed stays constant but the direction changes.}{Constant speed would require acceleration always perpendicular to velocity, which is not the case here because $a_x\neq 0$ affects the speed along the initial motion direction.}
  \Correct{Both the direction of motion and the speed change.}{The nonzero $a_y$ creates a growing $y$-velocity component, so the velocity direction rotates away from the $+x$ axis. Also, since $a_x\neq 0$, the component of acceleration parallel to the velocity is nonzero, so the speed (magnitude of $\vec v$) changes with time.}
}

\QBlock{Q9}{Two point charges are separated by a distance $r$ and exert an electrostatic force of magnitude $F$ on each other. If each charge's magnitude is reduced to one-third of its original value and the separation is doubled (to $2r$), what happens to the force magnitude?}{
  \Wrong{It stays the same.}{This would require the charge-product scaling to exactly match the $r^2$ scaling. Here both changes reduce the force, so it cannot remain unchanged.}
  \Wrong{It decreases by a factor of $9$.}{This accounts for $q_1q_2\to (1/3)(1/3)q_1q_2=\frac{1}{9}q_1q_2$ but incorrectly ignores the inverse-square dependence on distance.}
  \Correct{It decreases by a factor of $36$.}{By Coulomb's law, $F=k\frac{|q_1q_2|}{r^2}$. The charge product scales by $(1/3)(1/3)=1/9$ and the squared distance scales by $(2r)^2=4r^2$, so $F' = F\cdot \frac{1/9}{4}=\frac{F}{36}$; the force decreases by a factor of $36$.}
  \Wrong{It increases by a factor of $36$.}{This reverses the scaling in Coulomb's law. Reducing charge magnitudes and increasing separation both decrease the force rather than increase it.}
  \Wrong{It decreases by a factor of $4$.}{This includes the distance effect ($1/(2^2)=1/4$) but ignores that both charges were also reduced, which introduces an additional $1/9$ factor.}
}

\QBlock{Q10}{Escape speed from a spherical world is $v_e=\sqrt{\frac{2GM}{R}}$. World Y has mass $M$ and radius $R$. World X has mass $0.40M$ and radius $0.20R$. Which statement is correct?}{
  \Wrong{World X has a lower escape speed than World Y because it has smaller mass.}{This considers only $M$ and ignores the dependence on $R$. Escape speed depends on the ratio $M/R$, not on $M$ alone.}
  \Wrong{World X has the same escape speed as World Y because both $M$ and $R$ are smaller.}{Having both parameters smaller does not guarantee cancellation. You must compare $\frac{M}{R}$ for each world.}
  \Correct{World X has a higher escape speed than World Y (by a factor of about $1.41$).}{From $v_e\propto\sqrt{M/R}$, the ratio is $\frac{v_{e,X}}{v_{e,Y}}=\sqrt{\frac{0.40M/0.20R}{M/R}}=\sqrt{2}\approx 1.41$. Even though X is less massive, it is much more compact (smaller $R$), which increases $M/R$ and thus increases escape speed.}
  \Wrong{The escape speed depends on the spacecraft's mass, so it cannot be compared without that information.}{In Newtonian gravity (neglecting drag), $v_e=\sqrt{2GM/R}$ is independent of the escaping object's mass. The spacecraft mass cancels when deriving the expression from energy conservation.}
}

\QBlock{Q11}{A constant force $\vec F=(-2.5\ \mathrm{N})\,\hat\imath+(6.0\ \mathrm{N})\,\hat\jmath$ acts on a particle as it moves in the plane from $\vec r_i=(-1.2\,\hat\imath+0.5\,\hat\jmath)\ \mathrm{m}$ to $\vec r_f=(3.8\,\hat\imath+4.5\,\hat\jmath)\ \mathrm{m}$. What is the work done by the force?}{
  \Correct{$11.5\ \mathrm{J}$}{Work is the dot product $W=\vec F\cdot\Delta\vec r=F_x\Delta x+F_y\Delta y$. Here $\Delta\vec r=(3.8-(-1.2))\hat\imath+(4.5-0.5)\hat\jmath=5.0\hat\imath+4.0\hat\jmath$, so $W=(-2.5)(5.0)+(6.0)(4.0)=-12.5+24=11.5\ \mathrm{J}$.}
  \Wrong{$41.6\ \mathrm{J}$}{This comes from incorrectly using $W\approx |\vec F|\,|\Delta\vec r|$, which only works when force is parallel to displacement. Here the force is not parallel to the displacement.}
  \Wrong{$24.0\ \mathrm{J}$}{This uses only the $y$-component contribution $F_y\Delta y$ and ignores the negative work from the $x$-component. Work adds contributions from all components via the dot product.}
  \Wrong{$-12.5\ \mathrm{J}$}{This uses only $F_x\Delta x$ and ignores the positive work from the $y$-component. The total work is the sum $F_x\Delta x+F_y\Delta y$.}
}

\QBlock{Q12}{A $3.2\ \mathrm{kg}$ cart moving to the right at $4.5\ \mathrm{m/s}$ hits a bumper and rebounds to the left at $2.0\ \mathrm{m/s}$. The contact time is $0.15\ \mathrm{s}$. What is the magnitude of the average force on the cart during the collision?}{
  \Wrong{$96\ \mathrm{N}$}{This can result from using $\Delta v=4.5-2.0=2.5\ \mathrm{m/s}$ (forgetting the velocity reverses sign). The correct $\Delta v$ must use $v_f=-2.0\ \mathrm{m/s}$.}
  \Correct{$139\ \mathrm{N}$}{Use impulse-momentum: $F_{\text{avg}}=\frac{\Delta p}{\Delta t}=\frac{m(v_f-v_i)}{\Delta t}$. Taking right as positive, $v_i=+4.5\ \mathrm{m/s}$ and $v_f=-2.0\ \mathrm{m/s}$, so $\Delta v=-6.5\ \mathrm{m/s}$ and $|F_{\text{avg}}|=\left|\frac{3.2(-6.5)}{0.15}\right|\approx 1.39\times 10^2\ \mathrm{N}$.}
  \Wrong{$31\ \mathrm{N}$}{This is too small and typically comes from dividing by the wrong time interval or using only one speed instead of the change in velocity. Impulse depends on $v_f-v_i$.}
  \Wrong{None of the above.}{Option B matches the impulse calculation (to rounding). Therefore it is not correct that none of the above is correct.}
}

\QBlock{Q13}{A $72\ \mathrm{kg}$ snowboarder is initially at rest on nearly frictionless snow. A friend tosses a $250\ \mathrm{g}$ snowball directly toward the snowboarder at $90\ \mathrm{km/h}$, and the snowboarder catches it (the snowball sticks). What is the snowboarder's speed immediately after the catch?}{
  \Wrong{$0.31\ \mathrm{m/s}$}{This typically comes from forgetting to convert $90\ \mathrm{km/h}$ to $\mathrm{m/s}$ (treating $90$ as $90\ \mathrm{m/s}$). Unit conversion is essential before applying momentum conservation.}
  \Wrong{$25\ \mathrm{m/s}$}{This incorrectly assumes the snowboarder ends up moving at (or near) the snowball's initial speed. In a perfectly inelastic catch with a much larger combined mass, the final speed must be much smaller.}
  \Correct{$0.087\ \mathrm{m/s}$}{Treat the catch as perfectly inelastic and conserve momentum: $m_b v_b=(m_s+m_b)v_f$. Convert $90\ \mathrm{km/h}$ to $v_b=25\ \mathrm{m/s}$ and use $m_b=0.250\ \mathrm{kg}$, $m_s=72\ \mathrm{kg}$: $v_f=\frac{0.250\cdot 25}{72.25}\approx 0.0865\ \mathrm{m/s}\approx 0.087\ \mathrm{m/s}$.}
  \Wrong{$0.0087\ \mathrm{m/s}$}{This is a factor-of-10 error, often from a mistaken km/h conversion (e.g., dividing by $3.6$ twice) or misplacing a decimal in the final division.}
}

\end{document}
